\begin{frame}{PCA 목표와 5단계 절차}
\begin{center}\begin{varwidth}{0.92\textwidth}\usebeamercolor[fg]{normal text}
  \begin{itemize}
    \item \textbf{PCA 목표: 분산 최대화}
  \end{itemize}
  \vskip0.4em
  \begin{enumerate}
    
    \item \kb{공분산 행렬} $\Sigma = \frac{1}{m}X^{T}X$ 계산.
    \item $\Sigma$의 \kb{고유벡터}·\kb{고유값}을 구한다.
    \item 고유값 큰 순서대로 정렬 --- 각 고유벡터 = \kb{주성분},
      대응 고유값 = 그 방향의 분산 크기.
    \item 상위 $d$개의 고유벡터에 투영 $\to$ $d$차원 압축 표현.
  \end{enumerate}
\end{varwidth}\end{center}
\note{\begin{itemize}
\item 목표: 데이터를 저차원으로 \kb{투영}(project)하되,
      \textbf{투영된 데이터의 분산이 최대}인 방향을 찾는다
      = 정보를 가장 적게 잃는 방향.
\item 평균 0이 되도록 \kb{중심화}(centering).
\end{itemize}}
\end{frame}
